% 主旨：短训预算下稀疏奖励 PPO 学不会推塔，而采样后改写动作会破坏 PPO 比率；
%       本文把战术规则写成采样前、训练与推理共用的合法动作掩码，并给出三条贡献。
\section{引言}
\label{sec:intro}

多人在线战术竞技（multiplayer online battle arena, MOBA）游戏同时具备长时程决策、部分可观测和多分支动作三类难点，已成为深度强化学习的重要测试平台\cite{berner2019dota,ye2020full}。《王者荣耀》1v1 是其中规模较小、规则明确的任务：双方各控制一名英雄，在一条兵线上清兵、换血、补给并推进防御塔\cite{ye2020mastering,wei2022hok}。以鲁班七号为例，智能体每 \num{3} 帧需输出一个由按钮、移动方向、技能方向和目标单位组成的 \num{6} 分支动作，而每局可达 \num{20000} 帧。

已有的王者荣耀智能体依靠大规模分布式训练取得了职业水平。\textcite{ye2020mastering} 以近端策略优化（proximal policy optimization, PPO）\cite{ppo2017}为基础，引入目标注意力、长短期记忆网络（long short-term memory, LSTM）、基于游戏知识的动作掩码与双裁剪目标；\textcite{wei2022hok} 把该任务整理为公开环境并给出了基线。这类系统的训练以小时到天计。课程平台提供的 PPO 基线沿用了同一网络结构，但单次训练预算只有 \qtyrange{25}{60}{\minute}。

在这一预算下，基线面临一个探索问题。平台默认只给推塔奖励，未经约束的 \num{6} 分支随机探索几乎不会把英雄带到敌塔下并持续输出，外塔血量的变化主要来自小兵，与英雄动作关联很弱。官方基线训练 \qty{25}{\minute} 后，对内置规则对手的胜率为 \qty{0}{\percent}，敌方外塔保持满血（\cref{sec:ablation}）。直接的补救是把人工战术知识注入策略，但注入方式必须与策略梯度相容：若在策略采样之后用手写规则改写动作，例如低血时强制回城，改写后的动作不再来自策略分布，样本中记录的旧策略概率与实际执行的动作脱节，PPO 的重要性比率因此失去意义，分母接近零时还会数值不稳定。

针对这两点，在平台 PPO 基线上把鲁班七号的战术知识写成\emph{采样前}的合法动作掩码。规则只删除不应出现的动作，剩余动作仍由策略采样；同一掩码随样本保存，训练时用它计算新策略概率，与采样时记录的旧概率求比率，两者定义在同一动作支持集上。主要工作如下：
\begin{enumerate}
  \item 设计了训练与推理一致的战术动作掩码，覆盖射程门控、低血撤退与塔下避险、受压选敌、补给滞回和推进走廊五类规则。它用先验知识缩小策略的假设空间，而不引入特权信息；文中说明了它如何进入条件策略与 PPO 损失（\cref{sec:method}）。
  \item 设计了由 \num{11} 项组成的混合塑形奖励与按钮头冷启动偏置，为对线、补刀、生存和推塔提供稠密信号。
  \item 在腾讯开悟 V63.0.9 平台上完成三组同代码、同超参数的冷启动对比，并以续训模型的六局换边对战和离线引擎上的固定权重实验补充考察掩码的作用（\cref{sec:experiments}）。
\end{enumerate}

\cref{sec:related} 介绍相关工作，\cref{sec:problem} 给出任务设定，\cref{sec:method} 描述方法，\cref{sec:experiments} 报告实验，\cref{sec:conclusion} 总结并讨论局限。
